% =============================================================================
\section{Connections to Prior Theory}\label{sec:connections}
% =============================================================================

\subsection{From ERF to ECL}

Effective receptive field analyses show that gradients concentrate near the center with approximately Gaussian decay in deep convolutional networks~\citep{Luo2016EffectiveRF}. Our influence profile generalizes this to arbitrary perturbation kernels and embedding spaces. In linear regimes, \cref{prop:linear_influence} recovers squared impulse responses. The key advance is that ECL is defined for discrete, non-differentiable inputs (DNA sequences) and uses distribution-aware perturbations rather than gradient computation.

\subsection{Attribution methods as estimators of influence}

Integrated gradients~\citep{Sundararajan2017IG} and DeepLIFT~\citep{Shrikumar2017DeepLIFT} estimate feature contributions under differentiability and baseline-choice assumptions. Perturbation influence energies are complementary: they require only forward evaluations, apply to discrete inputs, and can be distribution-aware. More precisely, gradient methods estimate $\norm{\partial f/\partial x_i}^2$ (a local linearization at a single input), while perturbation methods estimate $\E[\norm{f(X) - f(X^{(i)})}^2]$ (a finite-difference average over $P_X$); these two estimands coincide in the linear regime by \cref{prop:linear_influence}, but diverge for nonlinear $f$ in ways that depend on the curvature of $f$ near typical inputs. ISM and its accelerations~\citep{Nair2022fastISM,Schreiber2022Yuzu} are computationally aligned with our influence definitions. A key distinction is that ISM at a reference sequence gives a \emph{point estimate} of local sensitivity, while our framework averages over the data distribution $P_X$, providing a population-level quantity.

\subsection{From RULER and NIAH to gNIAH}

The RULER benchmark~\citep{Hsieh2024RULER} defines ECL operationally as the maximum context length maintaining performance above a threshold. Our perturbation-variance ECL is more general: it applies at the embedding level without requiring a downstream task, is a continuous rather than binary measure, and provides statistical uncertainty quantification. The gNIAH protocol (\cref{sec:gniah}) adapts the Needle-in-a-Haystack concept to genomics by inserting known regulatory motifs at varying distances, providing a complementary probing approach.

\subsection{Representation comparison vs context utilization}

CKA~\citep{Kornblith2019CKA} measures similarity of representations across models or layers. ECL instead measures \emph{where representations come from} in input space. Both are essential for interpreting sequence models: CKA tells us whether two models learn similar features; ECL tells us whether they use similar input ranges to construct those features.

\subsection{Decision-theoretic ordering and information geometry}

Blackwell dominance (\cref{sec:blackwell}) provides a canonical comparison of context experiments. This connects to information geometry: the Fisher information matrix of the context-limited embedding distribution characterizes the local informativeness of the context experiment. Models with higher Fisher information at a given radius extract more useful signal from that context range.

\subsection{Connection to LongPPL and key-token identification}

LongPPL~\citep{Fang2024LongPPL} identifies ``key tokens'' whose predictions improve with longer context, achieving strong correlation with downstream performance. In our framework, key positions are those with large influence energy: $\cI(i;r) \gg 0$. Under the convention that $\mathrm{KeyScore}(i) = \log P(x_i \mid x_{<i}) - \log P(x_i \mid x_{i-w:i})$ is positive when the long context improves the prediction, LongPPL's difference-based identification can be seen as a special case of our out-of-window ablation $\Delta_{\mathrm{out}}$ applied to next-token prediction loss with a sign flip: $\mathrm{KeyScore}(i) \approx -\Delta_{\mathrm{out}}(w; i)$ under a log-loss discrepancy.
